\section{Determinization}
We use a type-directed transformation \(\Gamma \vdash e : \tau \darrow e'\) that preserves non-probabilistic structure and replaces probabilistic sampling by its expected value when the result is typed at expectation mode. Representative clauses are below; other constructs follow the obvious structural recursion.
\begin{mathpar}
  \inferrule[\textsc{Det-Uniform-E}]
    {\Gamma \vdash e_1 : \Float{\E} \darrow e_1' \\ \Gamma \vdash e_2 : \Float{\E} \darrow e_2'}
    {\Gamma \vdash \uniform(e_1,e_2) : \Float{\E} \darrow (e_1' + e_2') \times 0.5}
  \and
  \inferrule[\textsc{Det-Uniform-G}]
    {\Gamma \vdash e_1 : \Float{\G} \darrow e_1' \\ \Gamma \vdash e_2 : \Float{\G} \darrow e_2'}
    {\Gamma \vdash \uniform(e_1,e_2) : \Float{\G} \darrow \uniform(e_1',e_2')}
  \and
  \inferrule[\textsc{Det-Gauss-E}]
    {\Gamma \vdash e_1 : \Float{\E} \darrow e_1' \\ \Gamma \vdash e_2 : \Float{\G} \darrow e_2'}
    {\Gamma \vdash \gaussian(e_1,e_2) : \Float{\E} \darrow e_1'}
  \and
  \inferrule[\textsc{Det-Gauss-G}]
    {\Gamma \vdash e_1 : \Float{\G} \darrow e_1' \\ \Gamma \vdash e_2 : \Float{\G} \darrow e_2'}
    {\Gamma \vdash \gaussian(e_1,e_2) : \Float{\G} \darrow \gaussian(e_1',e_2')}
  \and
  \inferrule[\textsc{Det-Var}]
    {\Gamma \vdash x : \tau}
    {\Gamma \vdash x : \tau \darrow x}
  \and
  \inferrule[\textsc{Det-Abs}]
    {\Gamma, x{:}\tau_1 \vdash e : \tau_2 \darrow e'}
    {\Gamma \vdash \lambda x.\,e : \tau_1 \to \tau_2 \darrow \lambda x.\,e'}
  \and
  \inferrule[\textsc{Det-Rec}]
    {\Gamma, f{:}\tau_1 \to \tau_2, x{:}\tau_1 \vdash e : \tau_2 \darrow e'}
    {\Gamma \vdash \mathsf{rec}\;f\;x.\,e : \tau_1 \to \tau_2 \darrow \mathsf{rec}\;f\;x.\,e'}
  \and
  \inferrule[\textsc{Det-Const}]
    {\Gamma \vdash c : \Float{m}}
    {\Gamma \vdash c : \Float{m} \darrow c}
  \and
  \inferrule[\textsc{Det-Neg}]
    {\Gamma \vdash e : \Float{m} \darrow e'}
    {\Gamma \vdash -e : \Float{m} \darrow -e'}
  \and
  \inferrule[\textsc{Det-Add}]
    {\Gamma \vdash e_1 : \Float{m} \darrow e_1' \\ \Gamma \vdash e_2 : \Float{m} \darrow e_2'}
    {\Gamma \vdash e_1 + e_2 : \Float{m} \darrow e_1' + e_2'}
  \and
  \inferrule[\textsc{Det-Mul}]
    {\Gamma \vdash e_1 : \Float{\G} \darrow e_1' \\ \Gamma \vdash e_2 : \Float{\G} \darrow e_2'}
    {\Gamma \vdash e_1 \times e_2 : \Float{\G} \darrow e_1' \times e_2'}
  \and
  \inferrule[\textsc{Det-If}]
    {\Gamma \vdash e : \Bool \darrow e' \\ \Gamma \vdash e_1 : \tau \darrow e_1' \\ \Gamma \vdash e_2 : \tau \darrow e_2'}
    {\Gamma \vdash \ifkw\;e\;\thenkw\;e_1\;\elsekw\;e_2 : \tau \darrow \ifkw\;e'\;\thenkw\;e_1'\;\elsekw\;e_2'}
\end{mathpar}
Lambda, application, pairs, sums, and \(\letkw\) follow the same shape (recursing on subterms). The key point is that expectation-mode samples are determinized to their expected value, while general-mode samples remain probabilistic but have their subterms determinized.

